% !TeX program = xelatex
% Separate full model; not included in report.tex.
\documentclass[14pt]{extarticle}
\usepackage{fontspec}\usepackage[russian]{babel}
\setmainfont{DejaVu Sans}
\usepackage[paperwidth=84.1cm,paperheight=59.4cm,margin=1.5cm]{geometry}
\usepackage{tikz}\usetikzlibrary{calc,positioning}
\usepackage{array,booktabs}
\pagestyle{empty}\setlength{\parindent}{0pt}
\definecolor{ink}{HTML}{25364A}\definecolor{aux}{HTML}{6D7890}\definecolor{many}{HTML}{1D6F8B}
\tikzset{entity/.style={draw=ink,line width=1.2pt,rounded corners=5pt,fill=blue!4,text width=7.2cm,minimum height=3cm,align=center,font=\large,inner sep=9pt},rel/.style={draw=ink,line width=1.15pt,preaction={draw=white,line width=4pt}},service/.style={draw=aux,dashed,line width=1pt,preaction={draw=white,solid,line width=4pt}},mn/.style={draw=many,line width=2pt},tag/.style={fill=white,inner sep=3pt,font=\normalsize,text=ink}}
\begin{document}
\begin{center}
{\fontsize{32}{38}\selectfont\bfseries ЛабРезерв: полная инфологическая ER-модель}\par\vspace{5mm}
{\Large 15 сущностей на одном листе. Номера линий соответствуют перечню связей ниже.}\par\vspace{3mm}
{\large\textcolor{many}{\textbf{Синие линии — M:N}}\quad Сплошные линии — основные связи\quad\textcolor{aux}{Пунктир — авторство и исполнение действий}}
\end{center}
\begin{center}\begin{tikzpicture}[scale=.8,transform shape]

\node[entity] (role) at (0,11) {\textbf{Роль}\\Код: пользователь, техник, администратор};
\node[entity] (user) at (0,0) {\textbf{Пользователь}\\Идентификатор, логин, имя, активность};
\node[entity] (audit) at (0,-11) {\textbf{Запись аудита}\\Идентификатор, объект, действие, данные до/после, время};
\node[entity] (adm) at (13,11) {\textbf{Допуск}\\Идентификатор, срок действия, основание, дата отзыва};
\node[entity] (notice) at (13,0) {\textbf{Уведомление}\\Идентификатор, вид, текст, время создания и прочтения};
\node[entity] (conflict) at (13,-11) {\textbf{Отклонённый запрос}\\Идентификатор, запрошенный интервал, код отказа, время};
\node[entity] (type) at (26,11) {\textbf{Тип оборудования}\\Идентификатор, название, признак архивирования};
\node[entity] (book) at (26,0) {\textbf{Бронирование}\\Идентификатор, плановый интервал, цель, статус};
\node[entity] (session) at (26,-11) {\textbf{Фактический сеанс}\\Идентификатор, начало, окончание, комментарий};
\node[entity] (cap) at (39,11) {\textbf{Функциональная возможность}\\Идентификатор, название};
\node[entity] (alt) at (39,0) {\textbf{Предложенная альтернатива}\\Идентификатор, интервал, приоритет, состояние};
\node[entity] (repair) at (39,-11) {\textbf{Техническая работа / ремонт}\\Идентификатор, интервал, результат работы};
\node[entity] (equip) at (52,11) {\textbf{Единица оборудования}\\Идентификатор, инвентарный номер, название, характеристики};
\node[entity] (fault) at (52,0) {\textbf{Сообщение о неисправности}\\Идентификатор, описание, время, статус проверки};
\node[entity] (outage) at (52,-11) {\textbf{Период недоступности}\\Идентификатор, вид причины, интервал, причина, статус};
\draw[mn] (user.north) -- (role.south) node[tag,pos=.5]{R01: M:N};
\draw[rel] (user.north east) -- (adm.south west) node[tag,pos=.5]{R02};
\draw[rel] (adm.east) -- (type.west) node[tag,pos=.5]{R03};
\draw[mn] (type.east) -- (cap.west) node[tag,pos=.5]{R04: M:N};
\draw[rel] (type.north) -- (26,17) -- (52,17) -- (equip.north) node[tag,pos=.5]{R05};
\draw[rel] (user.south east) -- (5,-4) -- (26,-4) -- (book.south) node[tag,pos=.5]{R06};
\draw[rel] (equip.west) -- (46,11) -- (46,6) -- (26,6) -- (book.north) node[tag,pos=.5]{R07};
\draw[rel] (cap.south) -- (39,4) -- (30,4) -- (book.north east) node[tag,pos=.5]{R08};
\draw[rel] (book.east) -- (alt.west) node[tag,pos=.5]{R09};
\draw[rel] (equip.south west) -- (alt.north east) node[tag,pos=.5]{R10};
\draw[rel] (book.south west) -- (23,-5) -- (23,-8) -- (session.north west) node[tag,pos=.5]{R11};
\draw[rel] (user.east) -- (notice.west) node[tag,pos=.5]{R12};
\draw[rel] (book.west) -- (notice.east) node[tag,pos=.5]{R13};
\draw[rel] (user.north west) -- (-6,0) -- (-6,20) -- (59,20) -- (59,0) -- (fault.east) node[tag,pos=.5]{R14};
\draw[rel] (equip.south) -- (fault.north) node[tag,pos=.5]{R15};
\draw[rel] (fault.south) -- (outage.north) node[tag,pos=.5]{R16};
\draw[rel] (equip.east) -- (57,11) -- (57,-11) -- (outage.east) node[tag,pos=.5]{R17};
\draw[rel] (outage.west) -- (repair.east) node[tag,pos=.5]{R18};
\draw[rel] (user.south) -- (0,-5) -- (13,-5) -- (conflict.north) node[tag,pos=.5]{R19};
\draw[rel] (equip.north east) -- (61,14) -- (61,-18) -- (13,-18) -- (conflict.south) node[tag,pos=.5]{R20};
\draw[service] (user.north east) -- (6,2) -- (6,8) -- (adm.west) node[tag,pos=.5]{R21};
\draw[service] (user.west) -- (-7,-2) -- (-7,-21) -- (52,-21) -- (outage.south) node[tag,pos=.5]{R22};
\draw[service] (user.west) -- (-8,-3) -- (-8,-23) -- (55,-23) -- (outage.south east) node[tag,pos=.5]{R23};
\draw[service] (user.south west) -- (-5,-4) -- (-5,-16) -- (39,-16) -- (repair.south) node[tag,pos=.5]{R24};
\draw[service] (user.south) -- (audit.north) node[tag,pos=.5]{R25};
\end{tikzpicture}\end{center}
\vspace{2mm}
Кратности читаются у соответствующих сущностей: «Пользователь 1 — 0..N Бронирование» означает одного владельца каждой брони и любое число броней пользователя.
\vspace{5mm}

\begin{minipage}[t]{.49\linewidth}
\begin{tabular}{@{}lll@{}}\toprule № & Связь и смысл & Кратности\\\midrule
R01 & Пользователь — Роль: назначение ролей & 0..N — 0..N (M:N)\\
R02 & Пользователь — Допуск: получатель & 1 — 0..N\\
R03 & Допуск — Тип оборудования & 0..N — 1\\
R04 & Тип — Функциональная возможность & 0..N — 0..N (M:N)\\
R05 & Тип — Единица оборудования & 1 — 0..N\\
R06 & Пользователь — Бронирование: владелец & 1 — 0..N\\
R07 & Единица оборудования — Бронирование & 1 — 0..N\\
R08 & Функциональная возможность — Бронирование & 1 — 0..N\\
R09 & Бронирование — Предложенная альтернатива & 1 — 0..N\\
R10 & Единица оборудования — Альтернатива & 1 — 0..N\\
R11 & Бронирование — Фактический сеанс & 1 — 0..1\\
R12 & Пользователь — Уведомление: получатель & 1 — 0..N\\
R13 & Бронирование — Уведомление & 0..1 — 0..N\\
\bottomrule\end{tabular}\end{minipage}
\hfill
\begin{minipage}[t]{.49\linewidth}
\begin{tabular}{@{}lll@{}}\toprule № & Связь и смысл & Кратности\\\midrule
R14 & Пользователь — Сообщение: автор & 1 — 0..N\\
R15 & Единица оборудования — Сообщение & 1 — 0..N\\
R16 & Сообщение — Период недоступности & 0..1 — 0..1\\
R17 & Единица оборудования — Период недоступности & 1 — 0..N\\
R18 & Период недоступности — Ремонт & 1 — 0..N\\
R19 & Пользователь — Отклонённый запрос & 1 — 0..N\\
R20 & Единица оборудования — Отклонённый запрос & 1 — 0..N\\
R21 & Пользователь — Допуск: кто выдал & 1 — 0..N\\
R22 & Пользователь — Недоступность: кто зарегистрировал & 1 — 0..N\\
R23 & Пользователь — Недоступность: кто завершил & 0..1 — 0..N\\
R24 & Пользователь — Ремонт: исполнитель & 1 — 0..N\\
R25 & Пользователь — Запись аудита: автор действия & 0..1 — 0..N\\
\bottomrule\end{tabular}\end{minipage}
\vspace{5mm}
\textbf{Уточнения.} Для поломки R16 обязателен, для ТО сообщение отсутствует. Автор завершения R23 обязателен после восстановления.
Аудит хранит логическую ссылку на изменённый объект; отдельные FK ко всем сущностям не вводятся.
M:N раскрываются в даталогической модели таблицами user\_role и type\_capability; здесь они не являются отдельными концептуальными сущностями.
\end{document}
